% !TeX encoding = UTF-8
% !TEX root = main.tex
% Catena di compilazione attuale: latexmk -> .latexmkrc -> lualatex -> pdf

\listfiles
\documentclass[a4paper, italian, 11pt, oneside]{report}
\usepackage{z3f-preamble}% preambolo -> pacchetti custom -> pacchetti standard

%%%%%%%%%%%%%%%%%%%%
% NOTE DI UTILIZZO %
%%%%%%%%%%%%%%%%%%%%
% Nota 1: Aumentare Margini per rilegatura, bella l'efficienza nello spazio ma poco utile
% Nota 2: Trovare un'altra grafica per i bordi degli spazi di codice
% Nota 3: Cercare di rendere dinamico lo spazio di codice
% Nota 4: Prendere spunto dai template in giro per migliorare lo stile
% Nota 5: Riordinare i pacchetti in tematiche (ora che compilano)
% Nota 6: utilizzo di tectonic per la compilazione (rust)

%%%%%%%%%%%%
% SORGENTE %
%%%%%%%%%%%%

\begin{document}
\maketitle
%\frontmatter
\hypersetup{linkcolor=black}
\setcounter{tocdepth}{2}
\tableofcontents
%\elencoteoremi
% \cleardoublepage
% \phantomsection
% \addcontentsline{toc}{chapter}{\listfigurename}
% \listoffigures

% \cleardoublepage
% \phantomsection
% \addcontentsline{toc}{chapter}{\listtablename}
% \listoftables

% \addcontentsline{toc}{chapter}{\lstlistlistingname}
% \lstlistoflistings

%\mainmatter

% Battle arena
\chapter{Analisi}

% In questa sezione andrà fatta l’analisi dei requisiti e quella del problema, ossia verranno elencate le cose che l’applicazione dovrà fare (requisiti). Nell’analisi dei requisiti dell’applicazione si dovrà spiegare cosa l’applicazione dovrà fare. Non ci si deve concentrare sui particolari problemi, ma esclusivamente su cosa si desidera che l’applicazione faccia.

\section{Requisiti}

Il team di 4 developer si pone come obiettivo di realizzare un gioco 2D in stile 8bit, all’interno del quale un guerriero esplora un ambiente al fine di trovare 3 chiavi. Quest'ultime aprono il portale che identifica la fine del gioco con la vittoria del giocatore. 

Il guerriero può imbattersi in gruppi di nemici che pattugliano determinate aree e che deve sconfiggere per sopravvivere. Un gruppo può essere formato da 1 a N nemici. I nemici sono divisi in: 
\begin{itemize}
  \item Nemico ``base'' che, una volta sconfitto, ha la possibilità di lasciare a terra una pozione;
  \item Miniboss: un nemico con statistiche maggiorate(vita/danno/velocità...) a protezione dei forzieri delle chiavi e del portale. Una volta sconfitto, lascia a terra l'arma che ha equipaggiata;
\end{itemize}

I nemici si muovono in percorsi definiti e quando il guerriero entra nel loro raggio d'azione, procedono nella sua direzione.

Le chiavi si trovano all'interno di forzieri locati a loro volta dentro dei dungeon. I dungeon sono altri ambienti disaccoppiati dall'ambiente di partenza:
\begin{itemize}
  \item Caverna, accessibile tramite il suo ingresso;
  \item Casa, accessibile attraverso una porta;
  \item Labirinto, accessibile attraverso la sua entrata.
\end{itemize}

Il giocatore può incontrare 4 tipi di oggetti:
\begin{enumerate}
  \item Trasportabili:
  \begin{itemize}
    \item Consumabili: pozioni che possono essere lasciate a terra dai nemici base sconfitti e che recuperano la salute del giocatore;
    \item Di missione: chiavi;
    \item Equipaggiamento: scudo per parare i colpi e una spada forniti all'inizio del gioco. Altre armi che possono sostituire la spada iniziale del giocatore;
  \end{itemize}
  \item Non Trasportabili:
  \begin{itemize}
    \item Interagibili: 
    \begin{itemize}
      \item Forzieri presenti nella mappa che contengono altri oggetti;
      \item Portale per la vittoria.
    \end{itemize}
  \end{itemize} 
\end{enumerate}
Il guerriero può equipaggiare un'arma alla volta, lasciandola cadere in caso di raccolta di quella dei nemici ed un numero limitato di pozioni. Il guerriero può muoversi nelle 4 direzioni principali: su-giù, destra-sinistra.

Il giocatore ha accesso ad una interfaccia che riassume:
\begin{itemize}
  \item Salute disponibile;
  \item Pozioni raccolte;
  \item Chiavi trovate;
  \item Le statistiche dell'arma a disposizione.
\end{itemize}

Se la vita viene esaurita, il gioco finisce per game over. Il livello viene resettato al punto di partenza.
\newpage
Il giocatore può eseguire 4 azioni principali:
\begin{itemize}
  \item Attacco;
  \item Difesa;
  \item Interazione;
  \item Movimento.
\end{itemize}

\section{Modello del dominio}

% In questa sezione si descrive il modello del dominio applicativo, individuando le entità in gioco e in che modo sono in relazione le uno con le altre. Non vanno inserire idee circa possibili soluzioni implementative, ovvero sull’organizzazione interna del software. Infatti, la fase di analisi va effettuata prima del progetto: né il progetto né il software esistono nel momento in cui si effettua l’analisi. La discussione di aspetti propri del software (ossia, della soluzione al problema) appartengono alla sfera della progettazione, e vanno discussi successivamente. È obbligatorio fornire uno schema UML del dominio, che diventerà anche lo scheletro della parte di modello dell’applicazione, ovvero dei suoi elementi costitutivi (in ottica MVC - Model View Controller): se l’analisi è ben fatta, dovreste ottenere una gerarchia di concetti che rappresentano le entità che compongono il problema da risolvere. Un’analisi ben svolta prima di cimentarsi con lo sviluppo rappresenta un notevole aiuto per le fasi successive: è sufficiente descrivere a parole il dominio, quindi estrarre i sostantivi utilizzati, capire il loro ruolo all’interno del problema, le relazioni che intercorrono fra loro, e reificarli in interfacce.

Le entità che entrano in gioco sono: 

\chapter{Design}

\section{Architettura}

\section{Design dettagliato}

\chapter{Sviluppo}

\section{Testing automatizzato}

\section{Metodologia di lavoro}

\section{Note di sviluppo}

\end{document}
